% GENERATED FILE. Edit canonical lessons, then run npm run generate:books.
% canonical-source-hash: fnv1a64:aed41b48a2c84b20
% canonical-lessons: RU-C140-obedat, RU-C140-zavtrakat, RU-C140-uzhinat, RU-C140-stirat, RU-C140-gladit

\chapter{To have lunch, To have breakfast, To have dinner, To do the laundry, To iron}
\label{ch:ru-to-have-lunch-to-have-breakfast-to-have-dinner-to-do-the-laundry-to-iron}
\hlchaptermodality{140}

\begin{chapteropening}
\textbf{By the end of this chapter:} \emph{I can say to have lunch, to have breakfast, to have dinner, to do the laundry and to iron.}

Complete the last lesson of chapter 140: I can say to have lunch, to have breakfast, to have dinner, to do the laundry and to iron.
\end{chapteropening}

\section[\texorpdfstring{\ru{обедать} (obédat)}{обедать (obédat)}]{\ru{обедать} (obédat) — to have lunch}

\label{lesson:RU-C140-obedat}



\begin{quote}
Before the new one: say the Russian for an exit, then the Russian for a receipt.
\end{quote}

\subsection*{You'll want to know: \ru{обедать}}
\textbf{\ru{обедать}} (\emph{obédat}) — \textquotedblleft{}to have lunch\textquotedblright{}.

\begin{cousinweb}
\textbf{\ru{обедать}} is \emph{\ru{обед}}, \textquotedblleft{}lunch\textquotedblright{}, made into a verb with \emph{-\ru{ать}}. \emph{\ru{обед}} is \emph{\ru{об}-} with the root of \emph{\ru{есть}}, \textquotedblleft{}to eat\textquotedblright{} (Vasmer).
\end{cousinweb}

\begin{grammarlens}[title={What you've built}]
One more word for this chapter.
\end{grammarlens}

\subsection*{Your turn}
\begin{itemize}
\raggedright
  \item \emph{Say it:} \emph{obédat}
  \item \emph{Say it:} \emph{obédat}, once more
  \item \emph{Say it:} say \emph{obédat}
  \item \emph{Recall:} say the Russian for an exit, then the Russian for a receipt, then say \emph{obédat} again
\end{itemize}

\begin{tcolorbox}[breakable,colback=teal!4,colframe=teal!35!black,title={Before you move on}]
What does \ru{обедать} mean? (\textquotedblleft{}To have lunch\textquotedblright{}.) Say it once more.
\end{tcolorbox}
\section[\texorpdfstring{\ru{завтракать} (závtrakat)}{завтракать (závtrakat)}]{\ru{завтракать} (závtrakat) — to have breakfast}

\label{lesson:RU-C140-zavtrakat}



\begin{quote}
Before the new one: say the Russian for a receipt, then the Russian for to have lunch.
\end{quote}

\subsection*{You'll want to know: \ru{завтракать}}
\textbf{\ru{завтракать}} (\emph{závtrakat}) — \textquotedblleft{}to have breakfast\textquotedblright{}.

\begin{cousinweb}
\textbf{\ru{завтракать}} is \emph{\ru{завтрак}}, \textquotedblleft{}breakfast\textquotedblright{}, made into a verb; Old Russian already had \emph{\ru{завтрокать}} (Vasmer). \emph{\ru{завтрак}} is built on \emph{\ru{завтра}}, from \emph{\ru{за} \ru{утро}}, \textquotedblleft{}after the morning\textquotedblright{}.
\end{cousinweb}

\begin{grammarlens}[title={What you've built}]
One more word for this chapter.
\end{grammarlens}

\subsection*{Your turn}
\begin{itemize}
\raggedright
  \item \emph{Say it:} \emph{závtrakat}
  \item \emph{Say it:} \emph{závtrakat}, once more
  \item \emph{Say it:} say \emph{závtrakat}
  \item \emph{Recall:} say the Russian for a receipt, then the Russian for to have lunch, then say \emph{závtrakat} again
\end{itemize}

\begin{tcolorbox}[breakable,colback=teal!4,colframe=teal!35!black,title={Before you move on}]
What does \ru{завтракать} mean? (\textquotedblleft{}To have breakfast\textquotedblright{}.) Say it once more.
\end{tcolorbox}
\section[\texorpdfstring{\ru{ужинать} (úzhinat)}{ужинать (úzhinat)}]{\ru{ужинать} (úzhinat) — to have dinner}

\label{lesson:RU-C140-uzhinat}



\begin{quote}
Before the new one: say the Russian for to have lunch, then the Russian for to have breakfast.
\end{quote}

\subsection*{You'll want to know: \ru{ужинать}}
\textbf{\ru{ужинать}} (\emph{úzhinat}) — \textquotedblleft{}to have dinner\textquotedblright{}.

\begin{cousinweb}
\textbf{\ru{ужинать}} is \emph{\ru{ужин}}, \textquotedblleft{}supper\textquotedblright{}, made into a verb with \emph{-\ru{ать}}. \emph{\ru{ужин}} comes from Proto-Slavic \emph{*jugŭ}, \textquotedblleft{}south, noon\textquotedblright{}: first the midday meal (Vasmer).
\end{cousinweb}

\begin{grammarlens}[title={What you've built}]
One more word for this chapter.
\end{grammarlens}

\subsection*{Your turn}
\begin{itemize}
\raggedright
  \item \emph{Say it:} \emph{úzhinat}
  \item \emph{Say it:} \emph{úzhinat}, once more
  \item \emph{Say it:} say \emph{úzhinat}
  \item \emph{Recall:} say the Russian for to have lunch, then the Russian for to have breakfast, then say \emph{úzhinat} again
\end{itemize}

\begin{tcolorbox}[breakable,colback=teal!4,colframe=teal!35!black,title={Before you move on}]
What does \ru{ужинать} mean? (\textquotedblleft{}To have dinner\textquotedblright{}.) Say it once more.
\end{tcolorbox}
\section[\texorpdfstring{\ru{стирать} (stirát)}{стирать (stirát)}]{\ru{стирать} (stirát) — to do the laundry, to wash (clothes)}

\label{lesson:RU-C140-stirat}



\begin{quote}
Before the new one: say the Russian for to have breakfast, then the Russian for to have dinner.
\end{quote}

\subsection*{You'll want to know: \ru{стирать}}
\textbf{\ru{стирать}} (\emph{stirát}) — \textquotedblleft{}to do the laundry, to wash (clothes)\textquotedblright{}.

\begin{cousinweb}
\textbf{\ru{стирать}} comes from \emph{\ru{тирать}}, the repeated form of \emph{\ru{тереть}}, \textquotedblleft{}to rub\textquotedblright{} (Vasmer): clothes were washed by rubbing. Latin \emph{terō}, \textquotedblleft{}I rub\textquotedblright{}, is a cousin of \emph{\ru{тереть}}.
\end{cousinweb}

\begin{grammarlens}[title={What you've built}]
One more word for this chapter.
\end{grammarlens}

\subsection*{Your turn}
\begin{itemize}
\raggedright
  \item \emph{Say it:} \emph{stirát}
  \item \emph{Say it:} \emph{stirát}, once more
  \item \emph{Say it:} say \emph{stirát}
  \item \emph{Recall:} say the Russian for to have breakfast, then the Russian for to have dinner, then say \emph{stirát} again
\end{itemize}

\begin{tcolorbox}[breakable,colback=teal!4,colframe=teal!35!black,title={Before you move on}]
What does \ru{стирать} mean? (\textquotedblleft{}To do the laundry, to wash (clothes)\textquotedblright{}.) Say it once more.
\end{tcolorbox}
\section[\texorpdfstring{\ru{гладить} (gládit)}{гладить (gládit)}]{\ru{гладить} (gládit) — to iron; to stroke}

\label{lesson:RU-C140-gladit}



\begin{quote}
Before the new one: say the Russian for to have dinner, then the Russian for to do the laundry.
\end{quote}

\subsection*{You'll want to know: \ru{гладить}}
\textbf{\ru{гладить}} (\emph{gládit}) — \textquotedblleft{}to iron; to stroke\textquotedblright{}.

\begin{cousinweb}
\textbf{\ru{гладить}} $\leftarrow$ Proto-Slavic \emph{*gladiti}, \textquotedblleft{}to smooth\textquotedblright{} — inherited, with \emph{\ru{гладкий}}, \textquotedblleft{}smooth\textquotedblright{}. Vasmer relates it to Lithuanian \emph{glósti}, \textquotedblleft{}to stroke\textquotedblright{}, and Latvian \emph{glãstît}, \textquotedblleft{}to caress\textquotedblright{}.
\end{cousinweb}

\begin{grammarlens}[title={What you've built}]
One more word for this chapter.
\end{grammarlens}

\subsection*{Your turn}
\begin{itemize}
\raggedright
  \item \emph{Say it:} \emph{gládit}
  \item \emph{Say it:} \emph{gládit}, once more
  \item \emph{Say it:} say \emph{gládit}
  \item \emph{Recall:} say the Russian for to have dinner, then the Russian for to do the laundry, then say \emph{gládit} again
\end{itemize}

\begin{tcolorbox}[breakable,colback=teal!4,colframe=teal!35!black,title={Before you move on}]
What does \ru{гладить} mean? (\textquotedblleft{}To iron; to stroke\textquotedblright{}.) Say it once more.
\end{tcolorbox}
